\section{Representative Support Modes}

This section applies the scale-shape coordinates to several explicit support
registers. The examples are representatives only. They do not define a
canonical basis of the five-dimensional shape space.

Let
\[
\mathbf h_0=h_0\mathbf 1,
\qquad
h_0>0,
\]
be the reference regular support register. Every perturbed register below is
assumed to remain in the dodecahedral support chamber \(\mathcal C_D\).

\subsection{Regular state and uniform scale}

For the reference state
\[
\mathbf h=h_0\mathbf 1,
\]
one has
\[
\bar h=h_0,
\qquad
\mathbf r=\mathbf 0,
\qquad
\sigma=0,
\qquad
\eta=0.
\]

The realization is the regular dodecahedron:
\[
P(\mathbf h)=D.
\]

More generally, let
\[
\mathbf h=ch_0\mathbf 1,
\qquad
c>0.
\]

Then
\[
\bar h=ch_0,
\qquad
\mathbf r=\mathbf 0,
\qquad
\sigma=\eta=0,
\]
and
\[
P(\mathbf h)=cD.
\]

Thus the positive uniform ray changes only regular Euclidean scale.

\subsection{Balanced exchange between two face pairs}

Let
\[
\mathbf v_1=(1,-1,0,0,0,0)
\]
and consider
\[
\mathbf h
=
h_0\mathbf 1+a\mathbf v_1.
\]

Since \(\mathbf v_1\) has zero coordinate sum,
\[
\bar h=h_0.
\]

The residual is
\[
\mathbf r
=
a(1,-1,0,0,0,0),
\]
with magnitude
\[
\sigma
=
\sqrt{2}\,|a|.
\]

For \(a\ne0\),
\[
\widehat{\mathbf r}
=
\operatorname{sgn}(a)
\frac{1}{\sqrt{2}}
(1,-1,0,0,0,0),
\]
and
\[
\eta
=
\frac{\sqrt{2}\,|a|}{h_0}.
\]

One opposite-face pair moves outward relative to the mean, one moves inward,
and four remain unchanged. Since the residual is nonzero, the realization is
not homothetic to \(D\).

\subsection{One dominant face pair}

Let
\[
\mathbf v_2
=
(5,-1,-1,-1,-1,-1)
\]
and consider
\[
\mathbf h
=
h_0\mathbf 1+a\mathbf v_2.
\]

Again,
\[
\bar h=h_0.
\]

The residual satisfies
\[
\mathbf r
=
a(5,-1,-1,-1,-1,-1),
\]
and
\[
\sigma
=
\sqrt{30}\,|a|.
\]

For \(a\ne0\),
\[
\widehat{\mathbf r}
=
\operatorname{sgn}(a)
\frac{1}{\sqrt{30}}
(5,-1,-1,-1,-1,-1),
\]
with
\[
\eta
=
\frac{\sqrt{30}\,|a|}{h_0}.
\]

This register distinguishes one opposite-face pair from the remaining five.
A different dominant pair gives a symmetry-related example.

\subsection{A four-pair balanced mode}

Let
\[
\mathbf v_3
=
(1,1,-1,-1,0,0)
\]
and consider
\[
\mathbf h
=
h_0\mathbf 1+a\mathbf v_3.
\]

Then
\[
\bar h=h_0,
\qquad
\mathbf r
=
a(1,1,-1,-1,0,0),
\]
and
\[
\sigma=2|a|.
\]

For \(a\ne0\),
\[
\widehat{\mathbf r}
=
\operatorname{sgn}(a)
\frac{1}{2}
(1,1,-1,-1,0,0),
\]
while
\[
\eta
=
\frac{2|a|}{h_0}.
\]

Two face pairs move outward relative to the mean, two move inward, and two
remain at the mean support.

\subsection{Mixed scale and shape}

Consider
\[
\mathbf h
=
(h_0+\delta)\mathbf 1
+
a(1,-1,0,0,0,0),
\]
with
\[
h_0+\delta>0.
\]

Its mean support is
\[
\bar h=h_0+\delta,
\]
while its residual is
\[
\mathbf r
=
a(1,-1,0,0,0,0).
\]

Therefore
\[
\sigma
=
\sqrt{2}\,|a|
\]
and
\[
\eta
=
\frac{\sqrt{2}\,|a|}{h_0+\delta}.
\]

The same absolute residual has smaller relative magnitude when the mean
support is larger. This separates the roles of \(\sigma\) and \(\eta\).

\subsection{Normalized support registers}

For the balanced two-pair example,
\[
\mathbf h
=
h_0\mathbf 1+a(1,-1,0,0,0,0),
\]
the normalized register is
\[
\mathbf q
=
\frac{\mathbf h}{h_0}
=
\left(
1+\frac{a}{h_0},
1-\frac{a}{h_0},
1,1,1,1
\right).
\]

Its mean coordinate is one.

In general,
\[
\mathbf q
=
\frac{\mathbf h}{\bar h}
=
\mathbf 1+\eta\widehat{\mathbf r}.
\]

Thus normalized registers lie in the affine hyperplane
\[
\left\{
\mathbf q\in\mathbb R^6:
\frac{1}{6}\sum_i q_i=1
\right\}.
\]

The regular shape is represented by \(\mathbf 1\). Positive scalar multiples
of the same support register have the same normalized shape coordinate.

\subsection{Symmetry-related registers}

If \(g\) is a rotational symmetry and
\[
\mathbf h'=\rho(g)\mathbf h,
\]
then
\[
P(\mathbf h')=gP(\mathbf h).
\]

The two realizations are congruent, with permuted face-pair labels. They
satisfy
\[
\bar h(\mathbf h')=\bar h(\mathbf h),
\]
\[
\sigma(\mathbf h')=\sigma(\mathbf h),
\]
and
\[
\eta(\mathbf h')=\eta(\mathbf h).
\]

Their normalized residuals are related by
\[
\widehat{\mathbf r}(\mathbf h')
=
\rho(g)\widehat{\mathbf r}(\mathbf h).
\]

Examples in the same symmetry orbit therefore describe the same unlabeled
Euclidean shape.

\subsection{Admissibility}

The algebraic formulas are valid whenever the support coordinates are
positive. Their interpretation as dodecahedral modes requires
\[
\mathbf h\in\mathcal C_D.
\]

The local stability theorem guarantees this for sufficiently small
perturbations of \(h_0\mathbf 1\). The paper does not compute a global
maximum for \(a\) or \(\delta\).

At the boundary of \(\mathcal C_D\), a facet may become redundant, an edge
may contract, vertices may merge, or the incidence structure may change.
Such states lie outside the admitted dodecahedral mode family.

\subsection{Summary}

The examples display three basic cases:

\begin{enumerate}
\item uniform registers change only scale;
\item zero-mean registers change shape at fixed mean support; and
\item mixed registers change both scale and shape.
\end{enumerate}

For each register,
\[
\bar h
\]
records mean scale,
\[
\sigma
\]
records absolute shape departure,
\[
\eta
\]
records scale-free shape departure, and
\[
\widehat{\mathbf r}\in S^4
\]
records shape direction.

No mechanical or dynamical interpretation is introduced.
